\documentclass[11pt,oneside]{article}

\usepackage[a4paper,margin=27mm,headheight=15pt]{geometry}
\usepackage[T1]{fontenc}
\usepackage[utf8]{inputenc}
\IfFileExists{libertinus.sty}{\usepackage{libertinus}}{\usepackage{lmodern}}
\usepackage{microtype}
\usepackage{amsmath,amssymb,amsthm,mathtools,mathrsfs}
\usepackage{bm}
\usepackage{booktabs,longtable,array,tabularx}
\usepackage{graphicx}
\usepackage{xcolor}
\usepackage{enumitem}
\usepackage{fancyhdr}
\usepackage{titlesec}
\usepackage{listings}
\usepackage{xurl}
\usepackage{hyperref}
\usepackage[nameinlink,noabbrev]{cleveref}

\definecolor{linkblue}{RGB}{25,75,135}
\definecolor{shadegray}{RGB}{246,247,249}
\definecolor{rulegray}{RGB}{195,201,208}
\hypersetup{
  colorlinks=true,
  linkcolor=linkblue,
  citecolor=linkblue,
  urlcolor=linkblue,
  pdftitle={Sharp Decay Laws for the Fourier Transform of Rvachev's Up-Function},
  pdfsubject={Dyadic factorization, lacunary sine products, extremal envelopes, metric fluctuations, and Cesaro normalization},
  pdfkeywords={Rvachev up-function, Fabius function, Fourier transform, sinc product, lacunary product, doubling map, transfer operator, Thue-Morse}
}

\pagestyle{fancy}
\fancyhf{}
\fancyhead[L]{\small Fourier decay of Rvachev's up-function}
\fancyhead[R]{\small\nouppercase{\leftmark}}
\fancyfoot[C]{\thepage}
\renewcommand{\headrulewidth}{0.4pt}

\titleformat{\section}{\Large\bfseries}{\thesection}{0.7em}{}
\titleformat{\subsection}{\large\bfseries}{\thesubsection}{0.7em}{}
\titleformat{\subsubsection}{\normalsize\bfseries}{\thesubsubsection}{0.7em}{}
\setcounter{secnumdepth}{3}
\setcounter{tocdepth}{2}
\setlist[itemize]{leftmargin=2em,itemsep=0.25em,topsep=0.4em}
\setlist[enumerate]{leftmargin=2.3em,itemsep=0.25em,topsep=0.4em}
\setlength{\emergencystretch}{2em}

\newtheorem{theorem}{Theorem}[section]
\newtheorem{proposition}[theorem]{Proposition}
\newtheorem{lemma}[theorem]{Lemma}
\newtheorem{corollary}[theorem]{Corollary}
\newtheorem{conjecture}[theorem]{Conjecture}
\theoremstyle{definition}
\newtheorem{definition}[theorem]{Definition}
\newtheorem{example}[theorem]{Example}
\theoremstyle{remark}
\newtheorem{remark}[theorem]{Remark}
\newtheorem{warning}[theorem]{Warning}

\newcommand{\R}{\mathbb R}
\newcommand{\C}{\mathbb C}
\newcommand{\N}{\mathbb N}
\newcommand{\Z}{\mathbb Z}
\newcommand{\e}{\mathrm e}
\newcommand{\ii}{\mathrm i}
\newcommand{\dd}{\,\mathrm d}
\newcommand{\Up}{\operatorname{up}}
\newcommand{\sinc}{\operatorname{sinc}}
\newcommand{\Uhat}{\mathcal U}
\newcommand{\T}{\mathsf T}
\newcommand{\Lop}{\mathcal L}
\newcommand{\bigO}{\mathcal O}
\newcommand{\smallo}{o}
\newcommand{\esssup}{\operatorname*{ess\,sup}}
\newcommand{\vtwo}{\nu_2}
\newcommand{\dist}{\operatorname{dist}}

\newenvironment{boxedremark}{%
  \par\smallskip\noindent\begin{center}\begin{minipage}{0.94\textwidth}%
  \color{black}\hrule\vspace{0.6em}\small
}{%
  \vspace{0.6em}\hrule\end{minipage}\end{center}\smallskip
}

\lstdefinestyle{pseudo}{
  basicstyle=\ttfamily\small,
  backgroundcolor=\color{shadegray},
  frame=single,
  rulecolor=\color{rulegray},
  columns=fullflexible,
  keepspaces=true,
  showstringspaces=false,
  xleftmargin=0.7em,
  xrightmargin=0.7em,
  aboveskip=0.8em,
  belowskip=0.8em
}

\title{\bfseries Sharp Decay Laws for the Fourier Transform\\
of Rvachev's Up-Function\\[0.35em]
\large Dyadic factorization, lacunary fluctuations, extremal envelopes, and Ces\`aro normalization}
\author{Prepared for Vladimir Reshetnikov\\
\small Companion analysis for the \texttt{ProveIt} Fabius-function development}
\date{26 August 2026}

\begin{document}
\maketitle

\begin{abstract}
Let $\Up$ be Rvachev's compactly supported $C^\infty$ function, normalized to have
integral one, and use the Fourier convention
$\widehat h(\xi)=\int_{\R}h(t)\e^{-2\pi\ii\xi t}\dd t$.  Its Fourier transform is
\[
 \Uhat(\xi)=\prod_{n=0}^{\infty}
 \sinc\!\left(\frac{\pi\xi}{2^n}\right),
 \qquad \sinc z=\frac{\sin z}{z}.
\]
The common description ``roughly log-normal'' is correct only at the coarsest scale.
There is no single pointwise asymptotic equivalent: $\Uhat$ vanishes at every nonzero
integer, and after removal of the universal quadratic logarithmic term the remainder is
a lacunary Birkhoff sum for the doubling map.

This article gives an exact dyadic factorization and then separates four distinct notions
of amplitude.  The sharp global/extremal envelope has power exponent
\[
 \beta_\infty=\frac12+\frac{\log(2\pi/\sqrt3)}{\log2}
 =2.359014879111740\ldots .
\]
On almost every fixed dyadic ray $\xi=2^k y$, the pointwise exponent is
\[
 \beta_{\mathrm{typ}}=\frac12+\frac{\log(2\pi)}{\log2}
 =3.151496129472318\ldots ,
\]
with sharp positive law-of-the-iterated-logarithm excursions of order
$\sqrt{\log\xi\,\log\log\log\xi}$.  The absolute Ces\`aro mean is controlled by the
Fouvry--Mauduit constant $\lambda=0.661322602060565\ldots$ and has exponent
\[
 \beta_1=\frac12+\frac{\log(\pi/\lambda)}{\log2}
 =2.748070014871334\ldots .
\]
Finally, the exponent $\log(2\pi)/\log2$ appearing in the earlier draft is identified
exactly as the $L^2$ (root-mean-square) exponent, not a pointwise one.  A positive
real-analytic, eventually decreasing gauge is constructed that rigorously satisfies the
bounded-Ces\`aro requirement posed in the original question, and the residual dyadic
profile obstructing a naive all-$T$ limit is described by a transfer-operator eigenmeasure.
\end{abstract}

\tableofcontents
\newpage

\section{Statement of the problem and the main correction}

The object of study is
\begin{equation}\label{eq:def-U}
 \Uhat(x)=\prod_{n=0}^{\infty}
 \sinc\!\left(\frac{\pi x}{2^n}\right),
 \qquad x\in\R.
\end{equation}
It is real, even, entire, and rapidly decreasing on the real axis.  Since $\Up$ is
$C^\infty$ with compact support, integration by parts already gives
$\Uhat(x)=\bigO_N(|x|^{-N})$ for every $N$.  What that elementary argument does not
reveal is the much sharper scale of decay, nor the arithmetic oscillations hidden inside
that scale.

The earlier draft proposed, in effect, a deterministic equivalent of the form
\begin{equation}\label{eq:old-form}
 |\Uhat(x)|\stackrel{?}{\sim}
 C(\log x)^{-p}x^{-\log(2\pi)/\log2}
 \exp\!\left(-\frac{(\log x)^2}{2\log2}\right).
\end{equation}
There are three immediate structural difficulties with such a statement.

\begin{enumerate}
\item $\Uhat(m)=0$ for every nonzero integer $m$, so no positive pointwise equivalent can
hold on the full real axis.
\item Even on a fixed dyadic ray, the normalized remainder is a lacunary sine product.
For almost every ray it has unbounded, sharply quantifiable positive logarithmic excursions.
\item Pointwise, extremal, $L^1$, and $L^2$ amplitudes are governed by different
geometric means of the same lacunary product.  Conflating them produces a plausible but
incorrect exponent.
\end{enumerate}

The common factor in every correct notion of decay is
\begin{equation}\label{eq:universal-lognormal}
 \exp\!\left(-\frac{(\log x)^2}{2\log2}\right).
\end{equation}
The next power of $x$, however, depends on what ``amplitude'' means.  The following table
is a compact preview; all entries are proved below.  Write
\[
 G_\beta(x):=x^{-\beta}
 \exp\!\left(-\frac{(\log x)^2}{2\log2}\right),\qquad x>0.
\]

\begin{center}
\small
\begin{tabularx}{\textwidth}{@{}>{\raggedright\arraybackslash}p{0.24\textwidth}
 >{\centering\arraybackslash}p{0.16\textwidth}
 >{\centering\arraybackslash}p{0.21\textwidth}X@{}}
\toprule
Notion of size & effective sine mean $r$ & exponent
$\displaystyle\beta=\frac12+\frac{\log(\pi/r)}{\log2}$ & status of normalization \\
\midrule
Sharp global/extremal envelope
& $\sqrt3/2$
& $2.359014879\ldots$
& sharp up to $x^{o(1)}$; attained on a period-two dyadic ray \\
Root mean square ($L^2$)
& $1/\sqrt2$
& $\log(2\pi)/\log2=2.651496129\ldots$
& exact lacunary second moment \\
Absolute mean ($L^1$)
& $\lambda=0.661322602\ldots$
& $2.748070015\ldots$
& yields a bounded Ces\`aro ratio; dyadic-endpoint limit after scaling \\
Almost-everywhere pointwise
& $1/2$
& $3.151496129\ldots$
& with sharp one-sided LIL-scale lacunary excursions \\
\bottomrule
\end{tabularx}
\end{center}

\begin{boxedremark}
The exponent in the old draft is not meaningless.  It is exactly the $L^2$ exponent:
\[
 \frac12+\frac{\log(\pi\sqrt2)}{\log2}
 =\frac{\log(2\pi)}{\log2}.
\]
What fails is its interpretation as a pointwise decay law, together with the proposed
$(\log x)^{-1/2}$ correction.
\end{boxedremark}

\section{Fourier product, zeros, and the unique extrema}

\subsection{Normalization and product representations}

With the Fourier convention used in the abstract, the standard construction of the
up-function gives \eqref{eq:def-U}; see Arias de Reyna~\cite{arias1982} for a systematic
treatment.  The product converges normally on compact subsets of $\C$, since
$1-\sinc z=\bigO(z^2)$ near zero.

The elementary identity
\[
 \frac{\sin z}{z}=\prod_{m=1}^{\infty}
 \left(1-\frac{z^2}{\pi^2m^2}\right)
\]
yields a second, arithmetically revealing product.

\begin{proposition}[Integer zeros and their multiplicities]\label{prop:zeros}
For every $z\in\C$,
\begin{equation}\label{eq:weierstrass-collected}
 \Uhat(z)=\prod_{m=1}^{\infty}
 \left(1-\frac{z^2}{m^2}\right)^{1+\vtwo(m)},
\end{equation}
where $\vtwo(m)$ is the exponent of $2$ in $m$.  Thus the zero at each
$m\in\Z\setminus\{0\}$ has multiplicity $1+\vtwo(|m|)$.
\end{proposition}

\begin{proof}
Inserting the sine product into \eqref{eq:def-U} gives factors with zeros at
$z=2^nq$, $n\ge0$, $q\ge1$.  For a fixed positive integer $m$, the number of pairs
$(n,q)$ satisfying $m=2^nq$ is $1+\vtwo(m)$.  Since
$\sum_{m\ge1}(1+\vtwo(m))/m^2<\infty$, the collected genus-zero product converges
normally and rearrangement is legitimate.
\end{proof}

A cosine form, useful for comparison with the Fabius literature, follows from repeated
use of $\sinc(2z)=\sinc(z)\cos z$:
\begin{equation}\label{eq:cosine-form}
 \Uhat(z)=\prod_{m=1}^{\infty}
 \left(\cos\frac{\pi z}{2^m}\right)^m.
\end{equation}
The elementary scaling equation is
\begin{equation}\label{eq:scaling}
 \Uhat(x)=\sinc(\pi x)\Uhat(x/2),
 \qquad
 \Uhat(2x)=\sinc(2\pi x)\Uhat(x).
\end{equation}

\subsection{Exactly one extremum between consecutive zeros}

Although the present article is mainly about amplitudes, it is useful to close a small
logical gap from the original discussion: the oscillation has exactly one stationary
point in every interval between consecutive positive integers.

\begin{proposition}[Unique interlacing critical point]\label{prop:critical}
For each $n\ge1$, the restriction of $\Uhat$ to $(n,n+1)$ has exactly one critical
point.  It is the unique maximum or minimum on that interval.
\end{proposition}

\begin{proof}
Away from the integer zeros, logarithmic differentiation of
\eqref{eq:weierstrass-collected} gives
\begin{equation}\label{eq:log-derivative}
 \frac{\Uhat'(x)}{\Uhat(x)}
 =-2x\sum_{m=1}^{\infty}
 \frac{1+\vtwo(m)}{m^2-x^2}.
\end{equation}
The series and its derivative converge locally uniformly on
$\R\setminus\Z$.  Put
\[
 A(x):=\sum_{m=1}^{\infty}\frac{1+\vtwo(m)}{m^2-x^2}.
\]
For $x>0$ away from the integers,
\[
 A'(x)=2x\sum_{m=1}^{\infty}
 \frac{1+\vtwo(m)}{(m^2-x^2)^2}>0.
\]
On $(n,n+1)$ one has $A(x)\to-\infty$ as $x\downarrow n$ and
$A(x)\to+\infty$ as $x\uparrow n+1$.  Hence $A$ has exactly one zero there, and
\eqref{eq:log-derivative} gives exactly one critical point of $\Uhat$.  Since the
endpoints are zeros and the sign is constant in the open interval, that critical point
is the unique extremum.
\end{proof}

Let
\begin{equation}\label{eq:En}
 E_n:=\max_{n\le x\le n+1}|\Uhat(x)|.
\end{equation}
The sharp limsup law for $E_n$ will follow in \Cref{sec:extremal}.

\section{The exact dyadic factorization}\label{sec:dyadic}

Every $x\ge1$ has a unique representation
\begin{equation}\label{eq:mantissa}
 x=2^k y,\qquad k=\lfloor\log_2x\rfloor,\qquad 1\le y<2.
\end{equation}
Write
\[
 t:=y-1\in[0,1),\qquad a:=\log2,
\]
and define
\begin{align}
 P_n(t)&:=\prod_{j=0}^{n-1}|\sin(\pi2^jt)|,\label{eq:Pn}\\
 Q_n(t)&:=2^nP_n(t)
       =\prod_{j=0}^{n-1}|2\sin(\pi2^jt)|,\label{eq:Qn}\\
 H(y)&:=\prod_{m=1}^{\infty}
 \left|\sinc\!\left(\frac{\pi y}{2^m}\right)\right|.
 \label{eq:H}
\end{align}
The tail $H$ is continuous on $[1,2]$, strictly positive on $[1,2)$, and vanishes at
$y=2$ only because its first factor does.

\begin{theorem}[Exact dyadic factorization]\label{thm:factorization}
For $k\ge0$ and $1\le y<2$,
\begin{equation}\label{eq:exact-factorization}
 |\Uhat(2^ky)|
 =2^{-k(k+1)/2}(\pi y)^{-(k+1)}H(y)P_{k+1}(y-1).
\end{equation}
Equivalently,
\begin{equation}\label{eq:exact-factorization-Q}
 |\Uhat(2^ky)|
 =2^{-k(k+1)/2}(2\pi y)^{-(k+1)}H(y)Q_{k+1}(y-1).
\end{equation}
\end{theorem}

\begin{proof}
Split the product \eqref{eq:def-U} at $n=k$.  For $0\le n\le k$, set $j=k-n$.
Then
\[
 \left|\sin\!\left(\pi\frac{2^ky}{2^n}\right)\right|
 =|\sin(\pi2^jy)|
 =|\sin(\pi2^j(y-1))|.
\]
The product of the corresponding denominators is
\[
 \prod_{j=0}^{k}\pi2^jy
 =(\pi y)^{k+1}2^{k(k+1)/2}.
\]
The remaining factors, indexed by $n=k+m$ with $m\ge1$, are exactly $H(y)$.
This proves \eqref{eq:exact-factorization}; multiplying $P_{k+1}$ by
$2^{-(k+1)}Q_{k+1}$ gives \eqref{eq:exact-factorization-Q}.
\end{proof}

This identity isolates the entire problem.  The triangular denominator produces the
quadratic logarithm.  The mantissa $y$ contributes only bounded-scale factors.  All
nontrivial residual behavior is carried by the lacunary product $P_n(t)$.

\subsection{An exact logarithmic normal form}

For a real parameter $c$, define
\begin{equation}\label{eq:Phi}
 \Phi_c(b):=\frac{b^2}{2a}+\frac{bc}{a}-\frac b2-c.
\end{equation}
If $L=\log x=ka+b$, then direct expansion gives the elementary identity
\begin{equation}\label{eq:algebra-identity}
 -\frac a2k(k+1)-(k+1)(c+b)
 =-\frac{L^2}{2a}
  -\left(\frac12+\frac ca\right)L+\Phi_c(b).
\end{equation}
Taking $c=\log(2\pi)$ in \eqref{eq:exact-factorization-Q} gives the following exact,
not merely asymptotic, decomposition.

\begin{corollary}[Exact centered logarithm]\label{cor:centered-log}
Let $x=2^ky$, $1\le y<2$, $t=y-1$, and $L=\log x$.  Then
\begin{equation}\label{eq:exact-log-normal-form}
 \log|\Uhat(x)|
 =-\frac{L^2}{2a}-\beta_{\mathrm{typ}}L
  +\Psi_{\mathrm{typ}}(y)+S_{k+1}(t),
\end{equation}
where
\begin{align}
 \beta_{\mathrm{typ}}
   &:=\frac12+\frac{\log(2\pi)}{a},\label{eq:beta-typ-def}\\
 \Psi_{\mathrm{typ}}(y)
   &:=\log H(y)+\Phi_{\log(2\pi)}(\log y),\label{eq:Psi-typ}\\
 S_n(t)&:=\log Q_n(t)
 =\sum_{j=0}^{n-1}\log|2\sin(\pi2^jt)|.
 \label{eq:Sn}
\end{align}
At points where a sine factor vanishes, both sides are understood as $-\infty$.
\end{corollary}

\begin{boxedremark}
Formula \eqref{eq:exact-log-normal-form} is the decisive correction to a formal
Euler--Maclaurin treatment of the logarithm of the sinc product.  The singular function
$\log|\sin|$ cannot be replaced by a smooth average with a deterministic
$\log\log x$ remainder.  Its dyadic samples form the Birkhoff sum $S_n(t)$, whose
excursions are both arithmetically sensitive and, metrically, much larger than
$\log\log x$.
\end{boxedremark}

\section{Almost-everywhere pointwise decay and its sharp upper fluctuation}
\label{sec:typical}

Let $\T(t)=2t\pmod1$ be the doubling map and
\begin{equation}\label{eq:phi}
 \varphi(t):=\log|2\sin(\pi t)|.
\end{equation}
Then $S_n(t)=\sum_{j=0}^{n-1}\varphi(\T^jt)$.  The classical integral
\begin{equation}\label{eq:logsin-mean}
 \int_0^1\log|2\sin(\pi t)|\dd t=0
\end{equation}
shows that the Birkhoff sum is centered.

\begin{theorem}[Typical power exponent]\label{thm:typical}
For Lebesgue-almost every $t\in(0,1)$, and hence for almost every $y=1+t\in(1,2)$,
\begin{equation}\label{eq:Sn-sublinear}
 S_n(t)=o(n).
\end{equation}
Consequently, along the dyadic ray $x_k=2^ky$,
\begin{equation}\label{eq:typical-log-asymptotic}
 \log|\Uhat(x_k)|
 =-\frac{(\log x_k)^2}{2\log2}
  -\beta_{\mathrm{typ}}\log x_k+o(\log x_k),
\end{equation}
where
\begin{equation}\label{eq:beta-typ-value}
 \beta_{\mathrm{typ}}
 =\frac12+\frac{\log(2\pi)}{\log2}
 =3.151496129472318798\ldots .
\end{equation}
\end{theorem}

\begin{proof}
The function $\varphi$ belongs to $L^1(0,1)$ and the doubling map is ergodic for
Lebesgue measure.  Birkhoff's ergodic theorem and \eqref{eq:logsin-mean} give
\eqref{eq:Sn-sublinear}.  Substitute this into the exact identity
\eqref{eq:exact-log-normal-form}; the bounded mantissa term
$\Psi_{\mathrm{typ}}(y)$ is absorbed into $o(\log x_k)$.
\end{proof}

The $o(\log x)$ remainder is not remotely as small as a deterministic power of
$\log x$.  Aistleitner, Hofer, and Larcher proved a sharp one-sided law of the
iterated logarithm for the upper excursions of precisely this lacunary
product~\cite{ahl2017}.

\begin{theorem}[Sharp upper metric fluctuation]\label{thm:lil}
For Lebesgue-almost every $t\in(0,1)$,
\begin{equation}\label{eq:LIL-S}
 \limsup_{n\to\infty}
 \frac{S_n(t)}{\sqrt{n\log\log n}}=\pi.
\end{equation}
Equivalently, for almost every fixed $y\in(1,2)$, along $x_k=2^ky$,
\begin{multline}\label{eq:LIL-x}
 \limsup_{k\to\infty}
 \frac{
 \log|\Uhat(x_k)|
 +\dfrac{(\log x_k)^2}{2\log2}
 +\beta_{\mathrm{typ}}\log x_k
 -\Psi_{\mathrm{typ}}(y)
 }
 {\sqrt{\dfrac{\log x_k}{\log2}\,
 \log\log\log x_k}}
 =\pi.
\end{multline}
Changing $\log\log\log x_k$ by the asymptotically equivalent
$\log\log((\log x_k)/\log2)$ does not alter the limit.
\end{theorem}

\begin{proof}
Theorem~2 of Aistleitner--Hofer--Larcher states that, for every $\varepsilon>0$,
$Q_{L+1}(t)$ is at most
$\exp((\pi+\varepsilon)\sqrt{L\log\log L})$ for all sufficiently large $L$, and is
at least $\exp((\pi-\varepsilon)\sqrt{L\log\log L})$ for infinitely many $L$,
for almost every $t$.  Taking logarithms and shifting the index gives
\eqref{eq:LIL-S}.  Formula \eqref{eq:LIL-x} then follows from
\eqref{eq:exact-log-normal-form}, because $k+1\sim(\log x_k)/\log2$ and
$\log\log(k+1)\sim\log\log\log x_k$.
\end{proof}

\begin{remark}
The theorem says that the multiplicative correction to the typical core is eventually
at most
\[
 \exp\!\left((\pi+o(1))
 \sqrt{\frac{\log x}{\log2}\,\log\log\log x}\right),
\]
and reaches the same expression with $\pi-o(1)$ in place of $\pi+o(1)$ along a
subsequence of almost every dyadic ray.  These positive excursions are subpolynomial in
$x$ but larger than every fixed power of $\log x$.  No symmetric assertion about the
negative excursions is needed here; those are sensitive to exceptionally close returns to
the integer zero set.  In particular, no factor $(\log x)^c$, and especially no
$(\log x)^{-1/2}$ factor, can produce a pointwise constant-amplitude normalization.
\end{remark}

\begin{remark}[Exceptional rays]
The metric theorem does not erase arithmetic dependence.  If $t$ is dyadic rational,
then an iterate of $t$ is $0$ and the ray eventually lands on integer zeros.  If $t$ is
periodic under doubling, $P_n(t)$ has an exact periodic-geometric form.  By constructing
binary expansions with very long blocks of zeros, one can also force arbitrarily deep
negative excursions in $S_n(t)$.  Hence no nontrivial uniform lower envelope exists.
\end{remark}

\section{The sharp global and extremal envelope}\label{sec:extremal}

The largest possible long-term geometric mean of $P_n(t)$ is not $1$; the doubling
dynamics prevents an orbit from remaining near the maximum of $|\sin\pi t|$.  The exact
answer is attained by the period-two orbit $1/3\leftrightarrow2/3$.

\subsection{An ergodic-optimization lemma}

Set
\begin{equation}\label{eq:rho}
 \rho:=\frac{\sqrt3}{2}.
\end{equation}

\begin{lemma}[Maximum sine average]\label{lem:max-average}
Let $\mathcal M_\T$ denote the $\T$-invariant Borel probability measures on $[0,1]$.
Then
\begin{equation}\label{eq:max-average}
 \sup_{\mu\in\mathcal M_\T}
 \int_0^1\log|\sin(\pi t)|\dd\mu(t)
 =\log\rho.
\end{equation}
The unique maximizing measure is the equidistribution on the two-cycle
$\{1/3,2/3\}$.
\end{lemma}

\begin{proof}
Put $z=\sin^2(\pi t)$.  Under doubling, $z$ evolves by the logistic map
\[
 R(z)=4z(1-z),\qquad 0\le z\le1.
\]
Let $\eta$ be the push-forward of a $\T$-invariant measure.  If
$\int\log z\dd\eta=-\infty$, there is nothing to prove.  Otherwise invariance and
$\log R(z)=\log4+\log z+\log(1-z)$ imply
\begin{equation}\label{eq:log-one-minus-z}
 \int\log(1-z)\dd\eta(z)=-\log4.
\end{equation}
The elementary inequality
\begin{equation}\label{eq:z-inequality}
 z^3(1-z)\le\frac{27}{256},
\end{equation}
with equality only at $z=3/4$, gives
\[
 3\int\log z\dd\eta+\int\log(1-z)\dd\eta
 \le\log\frac{27}{256}.
\]
Using \eqref{eq:log-one-minus-z} yields
$\int\log z\dd\eta\le\log(3/4)$.  Since
$\log|\sin\pi t|=\tfrac12\log z$, this is \eqref{eq:max-average}.
Equality forces $z=3/4$ almost surely.  Its two lifts are $t=1/3,2/3$, exchanged by
$\T$, so invariance forces equal masses on the two points.
\end{proof}

\begin{lemma}[Maximum finite products]\label{lem:Mn}
Let $M_n:=\sup_{0\le t\le1}P_n(t)$.  Then
\begin{equation}\label{eq:Mn-limit}
 \lim_{n\to\infty}M_n^{1/n}=\rho.
\end{equation}
\end{lemma}

\begin{proof}
The cocycle identity
$P_{m+n}(t)=P_m(t)P_n(\T^mt)$ implies $M_{m+n}\le M_mM_n$, so
$\lim n^{-1}\log M_n$ exists by Fekete's lemma.  The standard empirical-measure
argument for an upper-semicontinuous potential gives
\[
 \lim_{n\to\infty}\frac1n\log M_n
 =\sup_{\mu\in\mathcal M_\T}
 \int\log|\sin(\pi t)|\dd\mu(t).
\]
For completeness, choose nearly maximizing $t_n$ and pass to a weak limit of the
empirical measures $n^{-1}\sum_{j<n}\delta_{\T^jt_n}$.  Every limit is invariant.
Truncating $\log|\sin\pi t|$ from below and applying upper semicontinuity gives the
upper bound.  For the reverse bound, ergodic decomposition reduces matters to an ergodic
invariant measure, and Birkhoff's theorem supplies a generic point.  Now apply
\Cref{lem:max-average}.
\end{proof}

\subsection{The envelope theorem}

Define
\begin{equation}\label{eq:beta-infty}
 \beta_\infty
 :=\frac12+\frac{\log(\pi/\rho)}{\log2}
 =\frac12+\frac{\log(2\pi/\sqrt3)}{\log2}
 =2.359014879111740707\ldots .
\end{equation}

\begin{theorem}[Sharp global logarithmic envelope]\label{thm:envelope}
One has
\begin{equation}\label{eq:global-limsup}
 \limsup_{x\to\infty}
 \frac{
 \log|\Uhat(x)|+\dfrac{(\log x)^2}{2\log2}}
 {\log x}
 =-\beta_\infty,
\end{equation}
where zeros contribute the value $-\infty$ and do not affect the limsup.
Equivalently,
\begin{equation}\label{eq:global-upper-o}
 |\Uhat(x)|\le
 \exp\!\left(-\frac{(\log x)^2}{2\log2}
 -(\beta_\infty-o(1))\log x\right)
\end{equation}
uniformly as $x\to\infty$, and the exponent cannot be improved.
\end{theorem}

\begin{proof}
In \eqref{eq:exact-factorization}, insert
$P_{k+1}(t)\le M_{k+1}=\rho^{k+1}\e^{o(k)}$ from
\Cref{lem:Mn}.  Apply \eqref{eq:algebra-identity} with
$c=\log(\pi/\rho)$.  Since $y\in[1,2)$ and $H(y)\le1$, the mantissa contribution is
bounded above, giving \eqref{eq:global-upper-o}.

For sharpness choose $y=4/3$, so $t=1/3$.  The doubling orbit alternates between
$1/3$ and $2/3$, and every sine factor equals $\rho$.  Thus
$P_{k+1}(1/3)=\rho^{k+1}$ exactly.  Formula \eqref{eq:exact-factorization} and
\eqref{eq:algebra-identity} then give
\begin{equation}\label{eq:exact-max-ray}
 |\Uhat(2^k\tfrac43)|
 =C_\star
 \exp\!\left(-\frac{(\log(2^k\tfrac43))^2}{2\log2}\right)
 (2^k\tfrac43)^{-\beta_\infty},
\end{equation}
with the positive constant
\[
 C_\star=H(\tfrac43)
 \exp\!\left(
 \Phi_{\log(\pi/\rho)}(\log\tfrac43)
 \right).
\]
This proves the matching lower limsup.
\end{proof}

The same exponent governs the largest local extrema.

\begin{corollary}[Sharp envelope of the stationary values]\label{cor:extrema-envelope}
For $E_n$ defined in \eqref{eq:En},
\begin{equation}\label{eq:En-limsup}
 \limsup_{n\to\infty}
 \frac{
 \log E_n+\dfrac{(\log n)^2}{2\log2}}
 {\log n}
 =-\beta_\infty.
\end{equation}
\end{corollary}

\begin{proof}
The upper bound follows from \Cref{thm:envelope}, uniformly for $x\in[n,n+1]$.
For the reverse bound, put $x_k=2^k(4/3)$ and $n_k=\lfloor x_k\rfloor$.  The unique
extremum in $[n_k,n_k+1]$ has magnitude at least $|\Uhat(x_k)|$.  Apply
\eqref{eq:exact-max-ray}; replacing $x_k$ by $n_k$ changes logarithms by $o(1)$.
\end{proof}

\begin{remark}
This is a sharp \emph{limsup envelope}, not an asymptotic formula for every $E_n$.
The sporadic rises visible in numerical plots are intrinsic: the binary orbit of the
mantissa can spend unusually long intervals near favorable or unfavorable parts of the
sine potential.
\end{remark}

\section{Absolute mean decay and a rigorous Ces\`aro gauge}\label{sec:cesaro}

The original question asked for a positive analytic, eventually monotone $g$ such that
\[
 \frac1T\int_{T_0}^{T}\frac{|\Uhat(x)|}{g(x)}\dd x
\]
remains between two positive constants, or ideally tends to one.  The bounded version
has a clean rigorous answer.

\subsection{The Fouvry--Mauduit constant}

Set
\begin{equation}\label{eq:I1n}
 I_1(n):=\int_0^1P_n(t)\dd t.
\end{equation}
Fouvry and Mauduit proved~\cite{fouvrymauduit1996} that there are constants
$\kappa>0$ and $0<\lambda<1$ such that
\begin{equation}\label{eq:FM}
 I_1(n)=\kappa\lambda^n(1+o(1)).
\end{equation}
Aistleitner, Hofer, and Larcher gave the rigorous enclosure
\begin{equation}\label{eq:lambda-rigorous}
 0.66130<\lambda<0.66135.
\end{equation}
A high-precision transfer-operator computation gives
\begin{equation}\label{eq:lambda-numeric}
 \lambda=0.661322602060565\ldots .
\end{equation}
Digits beyond \eqref{eq:lambda-rigorous} are numerical rather than claimed here as an
interval-certified proof.

The relevant transfer operator is
\begin{equation}\label{eq:transfer-L1}
 (\Lop f)(x)=\frac12\left[
 \sin\!\left(\frac{\pi x}{2}\right)f\!\left(\frac x2\right)
 +\cos\!\left(\frac{\pi x}{2}\right)f\!\left(\frac{x+1}{2}\right)
 \right],\qquad 0\le x\le1.
\end{equation}
Indeed, $I_1(n)=\int_0^1\Lop^n\mathbf1(x)\dd x$, and $\lambda$ is the dominant
eigenvalue of $\Lop$ on a suitable H\"older space.

\subsection{The \texorpdfstring{$L^1$}{L1} exponent and the gauge}

Define
\begin{equation}\label{eq:beta1}
 \beta_1:=\frac12+\frac{\log(\pi/\lambda)}{\log2}
 =2.748070014871334\ldots
\end{equation}
and
\begin{equation}\label{eq:g1}
 g_1(x):=x^{-\beta_1}
 \exp\!\left(-\frac{(\log x)^2}{2\log2}\right),\qquad x>0.
\end{equation}

\begin{theorem}[Bounded Ces\`aro normalization]\label{thm:cesaro-bounded}
There exist constants $0<A<B<\infty$ and $X_0>0$ such that, for every $T\ge2X_0$,
\begin{equation}\label{eq:cesaro-bounds}
 A\le\frac1T\int_{X_0}^{T}
 \frac{|\Uhat(x)|}{g_1(x)}\dd x\le B.
\end{equation}
Replacing $T$ by $T-X_0$ in the denominator does not change the conclusion.  Thus
$g_1$ rigorously satisfies the weaker requirement in the original question.
\end{theorem}

\begin{proof}
Let $x=2^ky$, $1\le y<2$, $t=y-1$, and $n=k+1$.  Apply
\eqref{eq:algebra-identity} with
$c=\log(\pi/\lambda)$.  The choice \eqref{eq:beta1} makes the powers of $x$ cancel,
and the exact factorization becomes
\begin{equation}\label{eq:ratio-exact-L1}
 \frac{|\Uhat(2^ky)|}{g_1(2^ky)}
 =W_1(y)\frac{P_n(t)}{\lambda^n},
\end{equation}
where
\begin{equation}\label{eq:W1}
 W_1(y):=H(y)
 \exp\!\left(\Phi_{\log(\pi/\lambda)}(\log y)\right)
 =H(y)\frac{\lambda}{\pi}
 y^{\beta_1-1}
 \exp\!\left(\frac{(\log y)^2}{2\log2}\right).
\end{equation}
The function $W_1$ is bounded on $[1,2]$ and has a positive minimum on
$[1,3/2]$.

Let
\[
 J_k:=\int_{2^k}^{2^{k+1}}
 \frac{|\Uhat(x)|}{g_1(x)}\dd x.
\]
After $x=2^ky$,
\begin{equation}\label{eq:Jk}
 J_k=2^k\lambda^{-n}
 \int_0^1W_1(1+t)P_n(t)\dd t.
\end{equation}
The symmetry $P_n(1-t)=P_n(t)$ gives
$\int_0^{1/2}P_n=I_1(n)/2$.  Hence the lower and upper bounds for $W_1$, together
with \eqref{eq:FM}, imply
\begin{equation}\label{eq:Jk-asymp-bounds}
 c\,2^k\le J_k\le C\,2^k
\end{equation}
for all sufficiently large $k$, with fixed $c,C>0$.

If $2^K\le T<2^{K+1}$, the sum of complete shells and the current partial shell is at
most a constant times $2^K$, hence at most a constant times $T$.  For the lower bound,
the complete preceding shell contributes at least $c2^{K-1}$, whereas
$T<2^{K+1}$.  This proves \eqref{eq:cesaro-bounds}.
\end{proof}

\subsection{Analyticity and monotonicity properties of the gauge}

\begin{proposition}[Eventual alternating derivatives]\label{prop:g-derivatives}
The function $g_1$ is positive and real analytic on $(0,\infty)$.  For each fixed
$m\ge0$, there exists $X_m$ such that
\begin{equation}\label{eq:derivative-sign}
 (-1)^m g_1^{(m)}(x)>0\qquad(x\ge X_m).
\end{equation}
In particular $g_1$ is eventually strictly decreasing, and every fixed derivative is
eventually monotone.
\end{proposition}

\begin{proof}
Put
\[
 u(x):=\frac{\log x}{\log2}+\beta_1.
\]
Then $g_1'(x)=-x^{-1}u(x)g_1(x)$.  Induction using
$\frac{\dd}{\dd x}=x^{-1}\frac{\dd}{\dd\log x}$ gives, for fixed $m$,
\begin{equation}\label{eq:derivative-asymp}
 \frac{g_1^{(m)}(x)}{g_1(x)}
 =x^{-m}\left((-1)^m u(x)^m+\bigO_m(u(x)^{m-1})\right).
\end{equation}
The leading term controls the sign once $x$ is large enough.
\end{proof}

The proposition gives eventual monotonicity separately for every derivative order.  It
does not assert complete monotonicity on one common half-line independent of $m$.

\subsection{Why a simple normalization need not converge for every \texorpdfstring{$T$}{T}}

The bounded result is elementary once \eqref{eq:FM} is known.  A finer spectral
statement explains the remaining obstruction to a single all-$T$ limit.

On a H\"older space, the operator \eqref{eq:transfer-L1} has a positive eigenfunction
$h$, a positive eigenmeasure $\nu$, and a spectral gap.  Normalize
$\Lop h=\lambda h$, $\Lop^*\nu=\lambda\nu$, and $\nu(h)=1$.  Then, for H\"older $q$,
\begin{equation}\label{eq:spectral-gap}
 \lambda^{-n}\Lop^nq
 =h\,\nu(q)+\bigO(\theta^n\|q\|_{\mathrm{Hol}})
 \qquad(0<\theta<1).
\end{equation}
It follows from \eqref{eq:Jk} that
\begin{equation}\label{eq:Jk-limit}
 \frac{J_k}{2^k}\longrightarrow C_1>0.
\end{equation}
After multiplying $g_1$ by $C_1$, the Ces\`aro average therefore tends to one along the
dyadic endpoints $T=2^K$.

For $T=2^Ky$ with fixed $1<y<2$, however, the newest partial shell has the same order as
all preceding shells together.  At continuity points of the limiting measure one obtains
a profile
\begin{equation}\label{eq:mantissa-profile}
 \lim_{K\to\infty}
 \frac1{2^Ky}\int_{1}^{2^Ky}
 \frac{|\Uhat(x)|}{C_1g_1(x)}\dd x
 =\frac{1+D(y)/C_1}{y},
\end{equation}
where
\begin{equation}\label{eq:D-profile}
 D(y)=\left(\int_0^1h(x)\dd x\right)
 \nu\!\left(W_1(1+\cdot)\mathbf1_{[0,y-1]}\right),
 \qquad D(2)=C_1.
\end{equation}
The right side of \eqref{eq:mantissa-profile} is a dyadic-mantissa profile rather than a
constant.

There is a conceptual reason.  The invariant equilibrium measure $\mu=h\nu$ for the
potential $\log|\sin\pi t|$ is singular.  Indeed, if it were absolutely continuous, the
uniqueness of the absolutely continuous invariant probability for the doubling map would
force it to be Lebesgue measure.  But Lebesgue measure gives the normalized pressure
value $-\log2$, whereas \eqref{eq:lambda-rigorous} gives
$\log\lambda>-\log2$.  Since $h$ is positive and continuous, $\nu$ is singular as well.
Thus the cumulative function in \eqref{eq:D-profile} is not linear.

The same argument applies if $g_1$ is multiplied by a positive continuous function of
$\{\log_2x\}$: this merely reweights the singular measure by a continuous density and
cannot turn it into Lebesgue measure.  Consequently, within this natural log-normal
power class (even allowing smooth log-periodic corrections), one should expect a
nonconstant dyadic profile rather than the full limit in the strongest version of the
original question.  This does not rule out deliberately scale-dependent, more exotic
gauges outside that class.

\section{A unified \texorpdfstring{$L^p$}{Lp} spectrum}\label{sec:Lp}

The preceding rates are instances of one pressure family.  For $p>0$, define
\begin{equation}\label{eq:Ip}
 I_p(n):=\int_0^1P_n(t)^p\dd t
\end{equation}
and the weighted transfer operator
\begin{equation}\label{eq:Lp-op}
 (\Lop_p f)(x):=\frac12\left[
 \sin^p\!\left(\frac{\pi x}{2}\right)f\!\left(\frac x2\right)
 +\cos^p\!\left(\frac{\pi x}{2}\right)f\!\left(\frac{x+1}{2}\right)
 \right].
\end{equation}
Let $\lambda_p$ be its dominant eigenvalue.  Then
\begin{equation}\label{eq:Ip-asymp}
 I_p(n)=\kappa_p\lambda_p^n(1+o(1)),
 \qquad \kappa_p>0,
\end{equation}
under the standard Perron--Frobenius theorem on a H\"older space.  Put
\begin{equation}\label{eq:rp-betap}
 r_p:=\lambda_p^{1/p},
 \qquad
 \beta_p:=\frac12+\frac{\log(\pi/r_p)}{\log2}.
\end{equation}

\begin{theorem}[$L^p$ shell normalization]\label{thm:Lp-shell}
For every fixed $p>0$, with
\[
 g_p(x):=x^{-\beta_p}
 \exp\!\left(-\frac{(\log x)^2}{2\log2}\right),
\]
there are constants $0<c_p<C_p<\infty$ such that, for all sufficiently large $k$,
\begin{equation}\label{eq:Lp-shell-bound}
 c_p2^k\le
 \int_{2^k}^{2^{k+1}}
 \left|\frac{\Uhat(x)}{g_p(x)}\right|^p\dd x
 \le C_p2^k.
\end{equation}
With the spectral asymptotic \eqref{eq:Ip-asymp}, the shell integral divided by $2^k$
has a positive finite limit after a constant rescaling of $g_p$.
\end{theorem}

\begin{proof}
Repeat the proof of \Cref{thm:cesaro-bounded}, replacing $\lambda$ by
$r_p=\lambda_p^{1/p}$ and raising the exact ratio to the $p$th power.  The bounded
mantissa factor is positive on $[1,3/2]$, and symmetry again supplies half of the
unweighted integral on that interval.
\end{proof}

\subsection{Three exact or limiting cases}

\paragraph{The $L^1$ case.}
Here $\lambda_1=\lambda$, so $\beta_1$ is exactly \eqref{eq:beta1}.

\paragraph{The $L^2$ case.}
There is an elementary identity:
\begin{equation}\label{eq:L2-exact}
 I_2(n)=2^{-n}.
\end{equation}
Indeed, $\Lop_2\mathbf1=\tfrac12\mathbf1$ because
$\sin^2u+\cos^2u=1$.  Hence $\lambda_2=1/2$,
$r_2=1/\sqrt2$, and
\begin{equation}\label{eq:beta2}
 \beta_2
 =\frac12+\frac{\log(\pi\sqrt2)}{\log2}
 =\frac{\log(2\pi)}{\log2}
 =2.651496129472318798\ldots .
\end{equation}
This is precisely the power in the earlier draft.

\paragraph{The limits $p\downarrow0$ and $p\to\infty$.}
The pressure variational principle is
\begin{equation}\label{eq:pressure}
 \log\lambda_p
 =\sup_{\mu\in\mathcal M_\T}
 \left(h_\mu(\T)-\log2
 +p\int\log|\sin\pi t|\dd\mu(t)\right).
\end{equation}
At $p=0$ the equilibrium state is Lebesgue measure, and differentiation of pressure gives
\[
 \lim_{p\downarrow0}r_p
 =\exp\!\left(\int_0^1\log|\sin\pi t|\dd t\right)=\frac12.
\]
Thus $\lim_{p\downarrow0}\beta_p=\beta_{\mathrm{typ}}$.  At zero temperature,
\eqref{eq:pressure} and \Cref{lem:max-average} give
\[
 \lim_{p\to\infty}r_p=\rho=\frac{\sqrt3}{2},
 \qquad
 \lim_{p\to\infty}\beta_p=\beta_\infty.
\]
The four headline rates are therefore not unrelated formulas; they are points on one
thermodynamic spectrum.

\section{Numerical evaluation of the mean constant}\label{sec:numerics}

The rigorous interval \eqref{eq:lambda-rigorous} is sufficient for every qualitative and
power-law conclusion above.  The extra digits in \eqref{eq:lambda-numeric} can be
reproduced very efficiently because the two inverse branches in \eqref{eq:transfer-L1}
contract by $1/2$.

Choose the Chebyshev--Lobatto nodes
\begin{equation}\label{eq:cheb-nodes}
 x_j=\frac{1-\cos(j\pi/N)}2,
 \qquad 0\le j\le N.
\end{equation}
Represent a function by its values at these nodes, evaluate it at $x_j/2$ and
$(x_j+1)/2$ using barycentric interpolation, and form the $(N+1)\times(N+1)$ matrix
corresponding to \eqref{eq:transfer-L1}.  Its Perron eigenvalue converges spectrally fast
in this problem.  A typical computation gives:

\begin{center}
\begin{tabular}{@{}rr@{}}
\toprule
$N$ & dominant collocation eigenvalue \\
\midrule
$8$  & $0.6613226020798928$ \\
$12$ & $0.6613226020605659$ \\
$16$ & $0.6613226020605659$ \\
$24$ & $0.6613226020605650$ \\
$32$ & $0.6613226020605658$ \\
\bottomrule
\end{tabular}
\end{center}

A compact implementation is conceptually as follows.

\begin{lstlisting}[style=pseudo,caption={Chebyshev collocation for the Perron eigenvalue}]
x[j] = (1 - cos(pi*j/N))/2,  j = 0,...,N
A = barycentric_interpolation_matrix(x, x/2)
B = barycentric_interpolation_matrix(x, (x+1)/2)
L = 0.5 * (diag(sin(pi*x/2))*A + diag(cos(pi*x/2))*B)
lambda_N = largest_real_eigenvalue(L)
\end{lstlisting}

For a fully certified decimal expansion, one may combine a high-precision approximate
eigenfunction with a Collatz--Wielandt enclosure:
if $v>0$ pointwise, then
\begin{equation}\label{eq:CW}
 \inf_x\frac{\Lop v(x)}{v(x)}
 \le\lambda\le
 \sup_x\frac{\Lop v(x)}{v(x)}.
\end{equation}
Interval arithmetic and a Chebyshev remainder bound can turn \eqref{eq:CW} into a
machine-checkable enclosure.

\section{Consequences and comparison with the earlier draft}\label{sec:comparison}

The rigorous picture can be summarized in a few precise corrections.

\subsection{The leading quadratic logarithm was right}

Every rate contains
\[
 -\frac{(\log x)^2}{2\log2}.
\]
It comes exactly from
$2^{-k(k+1)/2}$ with $k\sim\log_2x$.  Thus ``log-normal scale'' is an accurate
first description.

\subsection{The next power depends on the notion of amplitude}

There is no universal next coefficient.  The relevant exponents are
\begin{align*}
 \beta_\infty&=2.359014879111740\ldots
 &&\text{sharp global/extremal envelope},\\
 \beta_2&=2.651496129472318\ldots
 &&\text{root-mean-square amplitude},\\
 \beta_1&=2.748070014871334\ldots
 &&\text{absolute/Ces\`aro mean},\\
 \beta_{\mathrm{typ}}&=3.151496129472318\ldots
 &&\text{almost-everywhere pointwise amplitude}.
\end{align*}
The old coefficient $\log(2\pi)/\log2$ is exactly $\beta_2$.

\subsection{The missing correction is not \texorpdfstring{$\log\log x$}{log log x}}

For typical rays, the exact deterministic algebra supplies an additional
$-\tfrac12\log x$ relative to the old proposed pointwise coefficient.  After that, the
remainder is the lacunary sum $S_n(t)$, whose sharp metric size is
\[
 \sqrt{\log x\,\log\log\log x},
\]
not $\log\log x$.  Consequently a fixed power of $\log x$ cannot flatten the amplitude.

\subsection{Constant pointwise amplitude is impossible}

There are two independent obstructions:
\begin{itemize}
\item exact zeros at every nonzero integer;
\item unbounded positive lacunary excursions even after restriction to almost every dyadic ray.
\end{itemize}
The correct replacement is a hierarchy: sharp limsup envelope, metric pointwise law,
$L^p$ shell laws, and Ces\`aro bounds.

\subsection{The bounded Ces\`aro problem is solved}

The elementary analytic gauge
\[
 g_1(x)=x^{-2.748070014871334\ldots}
 \exp\!\left(-\frac{(\log x)^2}{2\log2}\right)
\]
satisfies the two-sided bound \eqref{eq:cesaro-bounds}; after constant normalization,
its averages converge to one at dyadic endpoints.  For arbitrary endpoints, a singular
dyadic-mantissa profile naturally remains.

\section{Formalization-ready decomposition}\label{sec:formalization}

The argument separates into modules of very different formal difficulty.  This is useful
for extending the existing Lean development.

\begin{enumerate}
\item \textbf{Elementary analytic module.}
Formalize normal convergence of \eqref{eq:def-U}, the scaling law
\eqref{eq:scaling}, the collected zero product \eqref{eq:weierstrass-collected}, and the
unique-critical-point proof \Cref{prop:critical}.

\item \textbf{Exact dyadic algebra.}
The factorization \Cref{thm:factorization} and identity
\eqref{eq:algebra-identity} are finite-product manipulations plus a tail product.  This
is the natural first target and yields an exact theorem with no asymptotic machinery.

\item \textbf{Ergodic optimization.}
The sharp envelope reduces to subadditivity, empirical measures, and the elementary
inequality $z^3(1-z)\le27/256$.  The only substantial library component is compactness of
probability measures and the variational passage for an upper-semicontinuous potential.

\item \textbf{Metric input.}
The Birkhoff theorem gives the typical exponent.  The sharp one-sided LIL constant $\pi$ should
initially be imported as an external theorem corresponding to Aistleitner--Hofer--Larcher,
rather than reproved inside the Fabius development.

\item \textbf{Mean input.}
Import the Fouvry--Mauduit asymptotic \eqref{eq:FM} as a named hypothesis/theorem.
Once available, the proof of \Cref{thm:cesaro-bounded} is elementary shell bookkeeping.
The exact $L^2$ identity \eqref{eq:L2-exact} is especially attractive for immediate
formalization.

\item \textbf{Transfer-operator refinement.}
The dyadic endpoint limit and singular profile require Perron--Frobenius theory on a
H\"older or bounded-variation space.  These can be isolated from the core decay theorems
and added later.
\end{enumerate}

\section{Conclusion}

The Fourier transform of Rvachev's up-function has a universal quadratic logarithmic
scale but no unique finer decay rate.  The exact factorization
\[
 |\Uhat(2^ky)|
 =2^{-k(k+1)/2}(\pi y)^{-(k+1)}H(y)
 \prod_{j=0}^{k}|\sin(\pi2^j(y-1))|
\]
shows why: after the deterministic triangular denominator is removed, one is left with a
lacunary multiplicative cocycle for the doubling map.

That cocycle has four natural summaries.  Its maximum ergodic mean is
$\sqrt3/2$, producing the sharp extremal exponent $\beta_\infty$.  Its almost-sure
geometric mean is $1/2$, producing $\beta_{\mathrm{typ}}$ together with sharp positive
iterated-logarithm excursions.  Its $L^1$ spectral radius is the Fouvry--Mauduit constant
$\lambda$, producing the Ces\`aro exponent $\beta_1$.  Its $L^2$ mean is exactly
$1/\sqrt2$, producing the exponent that appeared, without the correct interpretation,
in the earlier draft.

Thus the right answer is more structured than a single corrected equivalent.  The
transform decays on the log-normal scale
$\exp(- (\log x)^2/(2\log2))$, but its power correction records whether one asks for a
global envelope, a typical point, an absolute mean, or a root-mean-square mean.  Once
those notions are kept separate, the apparent irregularity is not noise: it is the
visible signature of dyadic dynamics.

\begin{thebibliography}{99}

\bibitem{arias1982}
J.~Arias de Reyna,
\emph{Definici\'on y estudio de una funci\'on indefinidamente diferenciable de soporte compacto},
Rev. Real Acad. Cienc. Exactas F\'is. Nat. (Madrid) \textbf{76} (1982), 21--38;
English translation, \emph{An infinitely differentiable function with compact support:
definition and properties}, arXiv:1702.05442 (2017),
\url{https://arxiv.org/abs/1702.05442}.

\bibitem{ahl2017}
C.~Aistleitner, R.~Hofer, and G.~Larcher,
\emph{On evil Kronecker sequences and lacunary trigonometric products},
Ann. Inst. Fourier (Grenoble) \textbf{67} (2017), no.~2, 637--687;
preprint title \emph{On parametric Thue--Morse sequences and lacunary trigonometric products},
arXiv:1502.06738,
\url{https://arxiv.org/abs/1502.06738}.

\bibitem{fouvrymauduit1996}
\'E.~Fouvry and C.~Mauduit,
\emph{Sommes des chiffres et nombres presque premiers},
Math. Ann. \textbf{305} (1996), 571--599.

\bibitem{baladi2000}
V.~Baladi,
\emph{Positive Transfer Operators and Decay of Correlations},
Advanced Series in Nonlinear Dynamics, vol.~16, World Scientific, 2000.

\bibitem{ruelle1978}
D.~Ruelle,
\emph{Thermodynamic Formalism: The Mathematical Structures of Classical Equilibrium
Statistical Mechanics},
Addison--Wesley, 1978; second edition, Cambridge University Press, 2004.

\bibitem{walters1982}
P.~Walters,
\emph{An Introduction to Ergodic Theory},
Graduate Texts in Mathematics, vol.~79, Springer, 1982.

\bibitem{reshetnikov-question}
V.~Reshetnikov,
\emph{Asymptotic decay rate of an infinite product of sinc functions},
Mathematics Stack Exchange question 2745497 (2018),
\url{https://math.stackexchange.com/questions/2745497/asymptotic-decay-rate-of-an-infinite-product-of-sinc-functions}.

\bibitem{proveit-draft}
V.~Reshetnikov,
\emph{Decay rate conjecture}, draft in the \texttt{ProveIt} repository,
\url{https://github.com/VladimirReshetnikov/ProveIt/blob/main/Analysis/FabiusFunction/docs/semi-formalized-research-frontiers/drafts/Decay%20rate%20conjecture/Decay%20rate%20conjecture.tex}.

\end{thebibliography}

\end{document}
